\subsection{Research Contribution}

The following six research contributions are presented as precise, technically grounded advances that map directly to the ESG ABSA research program described in prior sections. Each contribution is stated with clarity on novelty, methodological specificity, and the measurable criterion by which it will be evaluated. The focus remains on end-to-end ESG ABSA for bilingual Indonesian/English corporate disclosures, with explicit attention to the page-level, regulatory, and evidentiary needs of practitioners.

% \subsubsection{Contribution 1: An end-to-end OCR-to-record ESG processing pipeline from PDF to structured records}

% What this contribution delivers
% A fully integrated pipeline that ingests PDF sustainability reports and regulatory filings (including scanned and native PDFs), applies page- and layout-aware OCR, performs multi-language text extraction, normalizes and aligns content to a joint ESG taxonomy (GRI/SASB/TCFD), and outputs a structured, machine-readable record format (JSON/NER-ready) with sentence- and paragraph-level annotations and explicit ESG aspect mappings.
% Novelty
% Combines layout-aware OCR, multilingual text normalization, and cross-framework topic alignment into a single, reproducible data product specifically designed for ABSA in ESG disclosures. Unlike prior work that treats PDFs as plain text or bypasses layout, this pipeline preserves document structure to enable sentence-level ABSA and regulatory traceability. It integrates a sentence- and region-level provenance map that ties extracted content to exact document regions (sections, tables, footnotes) and to regulatory concepts, enabling auditable downstream analysis \cite{10.1109/icdar.2011.302,10.1093/jamiaopen/ooac045,10.21203/rs.3.rs-6644838/v1,10.3390/app122111063,10.1108/jd-03-2023-0055,10.1145/3151509.3151525,10.1145/3167132.3167236,10.1145/3322905.3322909}
% .
% The pipeline is designed for bilingual Indonesian/English reports, with integrated cross-language alignment modules and a joint ontology linking Indonesian terms to English counterparts and to GRI/SASB/TCFD anchors, facilitating cross-language comparability and regulatory consistency \cite{10.1016/j.inffus.2025.103073,10.1111/eufm.12509,10.48550/arxiv.2511.00024,10.3390/math12193119}

% Evaluation and metrics
% OCR quality metrics (character error rate, word error rate) stratified by language; layout accuracy metrics (section/column segmentation accuracy); sentence boundary detection precision/recall; ESG-topic alignment accuracy against a gold standard; and end-to-end regeneration fidelity (regression tests across 60–100 document samples).
% Deliverables
% A reusable, documented OCR-to-record toolkit, an explicit JSON schema for ESG aspects, language tags, section identifiers, and regulatory mappings; and a benchmark dataset with ground-truth region-level annotations for Indonesia-related disclosures.
% Rationale and grounding

% This contribution operationalizes the PDF-to-structure transformation essential for ABSA in ESG contexts. It builds on the demonstrated need for layout-aware processing and cross-language alignment in multilingual ESG texts, and leverages prior work on document structure extraction and regulatory mapping to deliver a robust, auditable data product \cite{10.1109/icdar.2011.302,10.1093/jamiaopen/ooac045,10.21203/rs.3.rs-6644838/v1,10.3390/app122111063,10.1108/jd-03-2023-0055,10.1145/3151509.3151525,10.1145/3167132.3167236,10.1145/3322905.3322909,10.21203/rs.3.rs-3600821/v1,10.2139/ssrn.3796152}. The Indonesian regulatory context (OJK/IDX) motivates an implementation that supports regulator-facing traceability and standardization \cite{10.22617/tcs210472-2,10.1108/mf-01-2023-0020,10.1007/s40804-021-00233-z}
% \subsubsection{Contribution 2: A page-aware, prompt-driven extraction framework with seven templates across bilingual reports}

% What this contribution delivers
% A page-aware extraction engine that uses seven deliberate prompt templates to drive extraction of ESG evidence, aspect terms, and sentiment signals from Indonesian and English pages. Templates are designed to respect page-level cues (headers, tables, footnotes, figures) and to capture evidence across multiple document regions (text blocks, tables, captions).
% Novelty
% Introduces a disciplined, multi-template prompt strategy for ABSA in ESG disclosures, explicitly engineered for bilingual Indonesian/English corpora. It extends beyond generic document-level prompts by injecting page-context signals, region identifiers, and regulatory anchors into prompts to improve alignment of extracted content with ESG aspects and regulatory concepts. This mirrors found work on prompt-template diversity and prompt ensembles but applies them to the ESG ABSA domain with cross-language considerations \cite{10.1609/aaai.v38i21.30574,10.48550/arxiv.2510.01222,10.18653/v1/2023.emnlp-main.975,10.1108/dta-01-2024-0065} .
% The seven templates are procedureally defined to elicit: (a) sentence-level aspect extraction, (b) sentence-level polarity per aspect, (c) cross-language term normalization, (d) regulatory mapping hints, (e) target/commitment vs. tone cues, (f) evidence tagging with document-region provenance, and (g) explanations of the rationale behind each extraction decision.
% Evaluation and metrics
% Template-wise ablation studies measuring extraction precision/recall for aspects and sentiments per template; inter-template agreement; and cross-language consistency of outputs. Stability metrics across prompts (variance in extractions for the same text region) will be reported.
% Deliverables
% The seven-template prompt suite, a registry of prompt configurations, and an evaluation report comparing template performance across Indonesian/English inputs, including guidance for practitioners on template selection and combination.
% Rationale and grounding

% This contribution directly addresses the practical need for robust extraction in ESG ABSA pipelines, especially in complex, bilingual reports where content spans paragraphs, tables, and captions. It aligns with literature on prompt-based NLP in finance/ESG contexts and the importance of domain-specific prompts for reliable outputs, while explicitly addressing cross-language inputs and regulatory alignment \cite{10.21203/rs.3.rs-6614220/v1,10.48550/arxiv.2510.01222,10.18653/v1/2023.emnlp-main.975,10.1108/dta-01-2024-0065,10.1609/aaai.v38i21.30574}

% \subsubsection{Contribution 3: A tone-aware ESG taxonomy distinguishing commitment, action, and outcome}

% What this contribution delivers
% A multi-dimensional taxonomy and annotation framework that separates tone from explicit commitments, actions, and outcomes. It operationalizes a three-tiered tone taxonomy (commitment, action, outcome) layered atop the ESG pillars (environment, social, governance) to enable fine-grained aspect-level ABSA that discriminates between intent, activity, and measurable results.
% Novelty
% Expands beyond conventional global sentiment or document tone to introduce a structured, action-oriented tonal ontology that maps to real-world governance and performance signals. This taxonomy enables researchers and practitioners to disentangle motivational language (commitment) from operational changes (action) and verifiable results (outcome), thereby reducing conflation and increasing interpretability for regulators and investors. This addresses a gap noted in ABSA literature regarding distinguishing tone from concrete sentiment toward ESG aspects \cite{10.1002/csr.70089,10.2139/ssrn.3738629,10.1108/mf-01-2023-0020}
% The taxonomy is designed for bilingual application, with cross-language alignment rules to ensure consistency of commitment/action/outcome labels across Indonesian and English expressions and across GRI/SASB/TCFD mappings.
% Evaluation and metrics
% Inter-annotator agreement on the taxonomy; cross-language agreement statistics; calibration of tone labels with actual disclosed outcomes (where targets or KPIs exist); precision/recall for each subcategory per aspect.
% Deliverables
% A formal, ontology-backed tone taxonomy, annotation guidelines, and an openly accessible reference corpus with tone-annotated sentences in both languages.
% Rationale and grounding

% This contribution directly addresses the need to distinguish tone from substantive ESG signals (commitment, action, outcome) to improve the fidelity and auditability of ABSA outputs. It resonates with cross-lingual ABSA and domain-adapted sentiment literatures calling for nuanced, topic- and context-sensitive sentiment representations and emphasizes interpretability for governance contexts \cite{10.1016/j.inffus.2025.103073,10.1111/eufm.12509,10.48550/arxiv.2511.00024,10.3390/math12193119,10.1002/csr.70089,10.2139/ssrn.3738629}

% \subsubsection{Contribution 4: An ontology-based ABSA layer for ESG aspect classification}

% What this contribution delivers
% An ontology-driven ABSA layer that provides explicit mappings from extracted phrases to ESG aspects (environmental, social, governance) and sub-aspects (e.g., emissions, energy efficiency, governance oversight, supply chain, diversity) with cross-reference to GRI/SASB/TCFD concepts. The layer includes an explicit hierarchical structure, relations among aspects, and an evidence-linking mechanism to regulatory sections.
% Novelty
% Introduces an explicit, machine-readable ESG ontology tailored to Indonesian/English disclosures, enabling consistent aspect classification across languages and frameworks. The ontology supports explainable ABSA by providing a pathway from textual cues to regulatory topics, metrics, and potential action items. While ABSA often relies on flat taxonomy or domain-adapted lexicons, this work elevates taxonomy to a formal ontology with hierarchical relationships and regulatory tag propagation.
% The ontology is designed to be extensible to additional frameworks (TCFD, ISSB) and to accommodate multilingual alignment through cross-lingual term mappings and synonym sets.
% Evaluation and metrics
% ontology coverage metrics (recall/precision of aspect mappings), hierarchical consistency checks, and cross-language mapping accuracy. End-user evaluation with regulator-facing scenarios (mapping to required disclosures) to validate interpretability.
% Deliverables
% The ESG ontology with a formal schema (OWL/RDF-compatible), mapping tables to GRI/SASB/TCFD indicators, and a bilingual synonym/translation file set for Indonesian-English terms.
% Rationale and grounding

% This contribution provides a principled backbone for ABSA in ESG analytics, enabling consistent, auditable aspect classification and regulatory traceability. It aligns with literature on ESG standardization and ontology-driven information extraction, addressing the need for structured, cross-framework alignment in a multilingual Indonesian context \cite{10.1108/medar-11-2020-1097,10.1108/aaaj-06-2021-5330,10.1108/aaaj-02-2020-4445,10.1111/1911-3838.12232,10.20900/jsr20210006}

% \subsubsection{Contribution 5: A ClimateBERT comparison framework for external semantic validation}

% What this contribution delivers
% A formal framework to compare ABSA-derived climate signals with external semantic validators (e.g., ClimateBERT-derived climate-target labels, ClimateBERT-NetZero outputs, and other domain-specific validators) to assess alignment and provide external validation of climate-related signals.
% Novelty
% Introduces a validation architecture that couples internal ABSA outputs with external climate-domain models to detect divergence, calibrate climate-specific signals, and identify when ClimateBERT-style labels provide additional predictive or regulatory value. This cross-model validation addresses concerns about relying solely on internal ABSA outputs for climate signaling and supports robust external verification for regulators and investors.
% The framework supports bilingual Indonesian/English data by enabling cross-language validation pathways, using multilingual or domain-adapted validators where appropriate.
% Evaluation and metrics
% Concordance measures (e.g., Cohen’s kappa) between ABSA-derived climate signals and ClimateBERT-derived classifications; predictive validity analyses linking approvals of climate targets to disclosed commitments; calibration analyses across languages.
% Deliverables
% A validation protocol, a computational framework for cross-model validation, and a set of guidance metrics for when to trust internal ABSA labels versus external validator signals.
% Rationale and grounding

% This contribution directly addresses the need for external semantic validation of climate-related ESG signals, leveraging established ClimateBERT frameworks and related domain models. It aligns with the literature on domain-specific climate NLP models and cross-model validation practices in ESG analytics \cite{10.21203/rs.3.rs-3600821/v1,10.48550/arxiv.2510.01222,10.18653/v1/2023.emnlp-main.975,10.1609/aaai.v38i16.29726,10.1108/dta-01-2024-0065,10.30519/ahtr.1436175,10.2139/ssrn.4308287}

% \subsubsection{Contribution 6: Explainability outputs including lexical triggers, TF-IDF coefficients, and ontology paths}

% What this contribution delivers
% A comprehensive explainability module that produces (i) lexical triggers (key terms and phrases that drive aspect-level sentiment), (ii) TF-IDF-based feature importance for per-aspect explanations, and (iii) ontology-path visualizations showing the reasoning chain from textual cues to ESG aspects and regulatory mappings.
% Novelty
% Combines rule-based lexical cues with data-driven feature importance (TF-IDF) and an ontology-path traversal to provide multi-faceted explanations for ABSA decisions. This triad supports regulator-friendly transparency by explicitly identifying why an extraction was made, what features influenced sentiment, and which ontology nodes guided the classification. This holistic explainability approach extends beyond standard post-hoc XAI methods (e.g., SHAP/LIME) by embedding ontology-informed reasoning into the explanation, improving traceability to ESG frameworks.
% The explainability outputs are designed to be language-agnostic at the architectural level, with bilingual support through lexicon and ontology mappings, facilitating cross-language interpretability and auditability.
% Evaluation and metrics
% Coverage of lexical triggers per aspect; correlation of TF-IDF weights with human judgments of importance; user studies with regulators to assess the clarity and usefulness of ontology-path visualizations; fidelity of explanations to the actual content in the PDF-to-record pipeline.
% Deliverables
% An explainability toolkit integrated into the ABSA pipeline, including: (a) a lexical trigger repository, (b) TF-IDF explanation dashboards, and (c) ontology-path visualization components.
% Rationale and grounding

% Explainability is central to ESG analytics for stakeholder trust and regulatory acceptance. This contribution synthesizes established XAI approaches with ontology-driven reasoning to deliver regulator-friendly outputs and actionable insights, addressing the imperative for interpretability highlighted across ABSA and ESG literature \cite{10.1002/csr.70089,10.2139/ssrn.3738629,10.1108/mf-01-2023-0020,10.3390/math12193119,10.2139/ssrn.3796152,10.1371/journal.pone.0288052,10.1051/e3sconf/202343602004}.
% Overall significance

% Each contribution is designed to be technically novel in the context of ESG ABSA for Indonesian/English bilingual disclosures, addressing practical regulatory needs (OJK/IDX) and the advanced analytical demands of investors. Collectively, they establish a comprehensive, end-to-end, auditable, and explainable ABSA framework that integrates document processing, sentence- and aspect-level analysis, tone semantics, regulatory alignment, cross-model validation, and user-centric explainability.
% Notes on references and grounding

% The proposed contributions draw on the references and themes discussed in prior sections, including the emergence of ClimateBERT and climate-target detectors, cross-lingual ABSA approaches, the need for sentence-level, aspect-level extraction, and the importance of explainability and reproducibility in ESG NLP pipelines. The references cited in the proposals align with ClimateBERT-related literature \cite{10.21203/rs.3.rs-3600821/v1,10.2139/ssrn.3796152,10.1111/1911-3846.12832,10.48550/arxiv.2510.01222,10.18653/v1/2023.emnlp-main.975,10.1609/aaai.v38i16.29726,10.1108/dta-01-2024-0065,10.30519/ahtr.1436175,10.2139/ssrn.4308287}, and cross-lingual ABSA work \cite{10.1016/j.inffus.2025.103073,10.1111/eufm.12509,10.48550/arxiv.2511.00024,10.3390/math12193119}, as well as ontology- and standard-focused ESG research \cite{10.1108/medar-11-2020-1097,10.1108/aaaj-06-2021-5330,10.1108/aaaj-02-2020-4445,10.1111/1911-3838.12232}

\subsubsection{Annotated bibliography (machine-readable XML)}

Note: The above references are included to illustrate the types of sources underpinning these contributions. In your finalized manuscript, please ensure that each claim and each methodological choice is supported by precise, citable sources with complete bibliographic details (DOIs, authors, venues, year). If you would like, I can expand the annotated bibliography with complete IEEE-formatted citations corresponding exactly to the six contributions and their evaluative criteria.